% !TEX root = Projektdokumentation.tex

% Hinweis: der Titel muss zum Inhalt des Projekts passen und den zentralen Inhalt des Projekts deutlich herausstellen
\newcommand{\titel}{Proxmox Backup Server}
\newcommand{\untertitel}{Einrichtung eines Proxmox Backup Servers und Integration in die ITA PVE-Umgebung}
\newcommand{\kompletterTitel}{\titel{} -- \untertitel}

\newcommand{\autorName}{Max Mustermann}
\newcommand{\autorAnschrift}{Musterstraße 1}
\newcommand{\autorOrt}{12345 Musterstadt}

\newcommand{\betriebLogo}{Firmen-Logo.png}
\newcommand{\betriebName}{Musterforschungsinstitut (MFI)}
\newcommand{\betriebAnschrift}{Beispielweg 1}
\newcommand{\betriebOrt}{12345 Musterstadt}

\newcommand{\ausbildungsberuf}{{\small Fachinformatiker (VO 2022) Fachrichtung: Systemintegration}}
\newcommand{\betreff}{Dokumentation zur betrieblichen Projektarbeit}
\newcommand{\pruefungstermin}{Sommer 2026}
\newcommand{\abgabeOrt}{Musterstadt}
\newcommand{\abgabeTermin}{26.05.2026}
